\documentclass[11pt]{amsart}
\usepackage{amsmath,amssymb,amsthm,mathtools,mathrsfs}
\usepackage[margin=1.08in]{geometry}
\usepackage{booktabs}
\usepackage{array}
\usepackage{longtable}
\usepackage{hyperref}
\hypersetup{hidelinks}
\usepackage[final,expansion=false]{microtype}
\setlength{\parskip}{4pt plus 1pt}
\emergencystretch=2.5em

\newtheorem{theorem}{Theorem}[section]
\newtheorem{proposition}[theorem]{Proposition}
\newtheorem{lemma}[theorem]{Lemma}
\newtheorem{corollary}[theorem]{Corollary}
\theoremstyle{definition}
\newtheorem{definition}[theorem]{Definition}
\newtheorem{problem}[theorem]{Problem}
\newtheorem{entry}[theorem]{Entry}
\theoremstyle{remark}
\newtheorem{remark}[theorem]{Remark}

\newcommand{\C}{\mathbb C}
\newcommand{\R}{\mathbb R}
\newcommand{\N}{\mathbb N}
\newcommand{\Z}{\mathbb Z}
\newcommand{\Q}{\mathbb Q}
\newcommand{\dd}{\,d}
\newcommand{\RH}{\mathrm{RH}}
\newcommand{\QW}{\mathcal Q_W}
\newcommand{\K}{\mathcal K}
\newcommand{\XiR}{\Xi}
\newcommand{\eps}{\varepsilon}
\newcommand{\EK}{\mathscr E_K}
\newcommand{\NK}{\mathcal N_K}
\DeclareMathOperator{\Spec}{Spec}
\DeclareMathOperator{\Tr}{Tr}
\DeclareMathOperator{\Var}{Var}
\DeclareMathOperator{\Ran}{Ran}
\DeclareMathOperator{\dist}{dist}

\newcommand{\NOGO}{\textup{[NO-GO]}}
\newcommand{\FALS}{\textup{[FALSIFIER]}}
\newcommand{\CTMD}{\textup{[COUNTERMODEL]}}
\newcommand{\SCAL}{\textup{[SCALE]}}

\title[Obstructions to Weil positivity]
{Obstructions to Weil Positivity:\\
A Catalogue of No-Go Theorems, Falsifiers\\
and Countermodels in the Prime--Gamma Doob Coordinate}

\author{David Alejandro Trejo Pizzo}
\thanks{Independent Researcher. \texttt{dtrejopizzo@gmail.com}}

\begin{document}

\begin{abstract}
The companion paper reduces the Riemann Hypothesis, along the
Connes--Consani--Moscovici and Suzuki route, to a single strict inequality
on a cofinal family of finite rows: the physical source surplus
\(\mathfrak g>\delta\).  This paper does not attempt that inequality.  It
does the complementary and, we argue, currently more useful thing: it
records in one place the mechanisms which provably cannot prove it.

We work throughout in the full-kernel Doob coordinate, in which Weil
positivity is the single sharp Poincar\'e inequality
\(\EK(r)\ge\frac12\Var_{\mu_K}(r)\), with \(\EK\) the positive jump energy
retaining every ordinary von Mangoldt atom and the complete Gamma jump
measure.  The catalogue contains thirty-one entries in eight families:
order and total positivity; coupling, transport and martingale routings;
local squares, sums of squares and Hodge--Korn completions; determinants
and minor charging; finite heads and adaptive filtering; spectral and
operator-theoretic detectors; realization and factorization rigidity; and
the semilocal branch-selection coordinate.  Each entry is classified as a
no-go theorem inside the literal Riemann system, a conditional falsifier
valid under the off-line counterfactual, a countermodel satisfying all
structural hypotheses of a proposed mechanism, or a scale obstruction.

Three structural principles account for every entry.  \emph{Radical
saturation}: the infinite-dimensional family \(K^{(2j)}/K\) attains
equality, and its first member is unbounded, so no decomposition into
nonnegative summands may contain a local edge square.  \emph{Measure
capacity}: the source carries singular mass on the lines
\(t\pm s=\log p^k\) while the polar pair law is absolutely continuous, so
no endpoint-preserving positive coupling exists.  \emph{Counterfactual
invariance}: under the existence of an off-line orbit, every positive
datum of the system persists while the strict surplus fails, so no
implication insensitive to the location of the mean-periodic frequencies
can succeed.

We also record the unconditional structure that survives, including an
exact and fully explicit Gamma absorption theorem
\(2K(u)<\frac{1501}{2000}r_\Gamma(u)\) which removes the Hermitian
connection power with a fixed margin, and a rigidity theorem showing that
the remaining contraction threshold is invariant under every change of
signed realization.  We conclude by delimiting the narrow corridor of
mechanisms that remain admissible.  No statement in this paper proves the
Riemann Hypothesis, and by the counterfactual invariance principle none
of them could.
\end{abstract}

\maketitle
\tableofcontents

\section{Purpose}

A research programme that reduces a famous conjecture to a single
inequality produces two kinds of output.  The first is the reduction
itself.  The second, usually discarded, is the record of everything that
was tried against the surviving inequality and provably cannot work.

The second output is the subject of this paper.  Our reason for
publishing it separately is stated precisely in
Section~\ref{sec:circularity}: in the coordinate developed here, the
surviving inequality is not a preliminary estimate from which Weil
positivity follows.  It \emph{is} Weil positivity, evaluated at a
stationary residual.  Consequently a catalogue of obstructions is not a
consolation prize for having failed to close the gate.  It is the part of
the work whose validity does not depend on the gate ever being closed,
and it is directly useful to anyone else approaching the same wall.

We have tried to make each entry falsifiable in the ordinary sense: it
either exhibits an explicit object in the literal Riemann system, or it
exhibits a model satisfying a stated list of hypotheses and violating the
proposed conclusion.  Where an entry rules out a mechanism we ourselves
proposed in earlier work, we say so.

\subsection{What is not claimed}

We do not claim that the Riemann Hypothesis is false, nor that the
Connes--CCM route is unviable.  Every entry below rules out a
\emph{method}, not the conclusion.  We also do not claim the catalogue is
exhaustive; it is exhaustive only over the mechanisms we tested.

\section{The canonical coordinate}
\label{sec:coordinate}

All obstructions are stated in one coordinate, so that they can be
compared.  We recall its construction.

Let \(\xi(s)=\frac12s(s-1)\pi^{-s/2}\Gamma(s/2)\zeta(s)\) and
\(\XiR(z)=\xi(\frac12+iz)\).  Let \(K\) be Riemann's even kernel,
normalized by
\[
 \widehat K=\XiR .
\tag{2.1}
\]
The theta representation gives, for \(x\ge0\),
\[
 K(x)=2\pi e^{5x/2}\sum_{m\ge1}m^2\bigl(2\pi m^2e^{2x}-3\bigr)
 e^{-\pi m^2e^{2x}},
\tag{2.2}
\]
in which every summand is positive because \(2\pi m^2e^{2x}>3\).  Hence
\(K\) is smooth, even, strictly positive and double-exponentially
decreasing.  Put
\[
 h(x)=\cosh(x/2),
 \qquad
 c_K=\int_\R h(x)K(x)\dd x=\widehat K(i/2)=\xi(0)=\tfrac12 .
\tag{2.3}
\]

\begin{theorem}[Full-kernel Doob--variance identity]
\label{thm:doob}
For every even compactly supported smooth multiplier \(r\),
\[
 \boxed{
 \QW(Kr,Kr)=\EK(r)-\tfrac12\Var_{\mu_K}(r),}
\tag{2.4}
\]
where
\[
 \EK(r)=\int_0^\infty\!\!\int_\R
 K(x)K(x-u)|r(x)-r(x-u)|^2\dd x\dd\nu_\zeta(u),
\tag{2.5}
\]
\[
 \dd\nu_\zeta(u)=\sum_{n\ge2}\frac{\Lambda(n)}{\sqrt n}\delta_{\log n}(\dd u)
 +\frac{e^{-u/2}}{1-e^{-2u}}\dd u,
 \qquad
 \dd\mu_K=\frac{hK\dd x}{c_K}.
\tag{2.6}
\]
The identity is unconditional.  Consequently the Riemann Hypothesis is
equivalent to the single sharp Poincar\'e inequality
\[
 \boxed{\EK(r)\ge\tfrac12\Var_{\mu_K}(r)}
\tag{2.7}
\]
for all even compactly supported smooth \(r\).
\end{theorem}

The virtue of (2.4) is that it is an identity, not an estimate: no prime,
Gamma or polar term has been separated, bounded, or discarded.  Every
obstruction below is therefore an obstruction to proving (2.7), and can
be compared with every other.  We write \(\mathcal J_u(r)\) for the inner
displacement form in (2.5),
\[
 \mathcal J_u(r)=\int_\R K(x)K(x-u)|r(x)-r(x-u)|^2\dd x\ \ge0 ,
\tag{2.8}
\]
so that \(\EK(r)=\int_0^\infty\mathcal J_u(r)\dd\nu_\zeta(u)\).

\subsection{The radical}

Let \(L\ge0\) be the self-adjoint generator associated with the closed
form \(\EK\) on \(L^2_{\rm even}(\mu_K)\), and set
\[
 \mathcal R=\overline{\operatorname{span}}
 \Bigl\{\,\tfrac{K^{(2j)}}{K}\ :\ j\ge1\,\Bigr\}_{\rm centered},
 \qquad
 \mathscr C=(\mathbf 1\oplus\mathcal R)^\perp .
\tag{2.9}
\]
Then \(L\mathbf 1=0\) and \(Lr=\frac12r\) for \(r\in\mathcal R\): the
radical is exactly the threshold eigenspace.  Equivalently, in multiplier
coordinates, \(\mathcal R^\perp\) is characterized by the convolution
equation
\[
 (hq)*K=0 ,
\tag{2.10}
\]
whose analytic solutions are mean-periodic with frequencies in the zero
divisor of \(\XiR\), including the explicit modes
\(\chi_z(x)=\cos(zx)/h(x)\).  We write
\(\NK=\{F:F*K=0\}\).

\section{Unconditional structure that survives}
\label{sec:positive}

Before the catalogue we record what is proved, since several obstructions
are sharp precisely because these positive facts are available.

\begin{theorem}[Essential threshold]
\label{thm:ess}
\(\sigma_{\rm ess}(L|_{\mathbf 1^\perp})\subset[1/2,\infty)\).  Hence,
after shorting the exact threshold radical, every spectral point of
\(L|_{\mathscr C}\) in \((0,1/2)\) is an isolated eigenvalue of finite
multiplicity.
\end{theorem}

The proof combines local compactness of the form domain, obtained from the
Gamma small-jump estimate
\(\mathscr E_\Gamma+C\|f\|^2\gtrsim\int\log(2+|\xi|)|\widehat f(\xi)|^2\dd\xi\),
with a moving prime-number-theorem quadrature giving the exact tail mass
\(2c_K^2=1/2\).

\begin{theorem}[Heat-core Galerkin exhaustion]
\label{thm:heat}
The form \(\EK\) is closed on its maximal domain.  With
\(S_{\mathscr C}=L|_{\mathscr C}+\frac12I\ge\frac12I\) and any
Hilbert-dense \((g_j)\subset\mathscr C\), the nested spaces
\(V_M^{\rm heat}=\operatorname{span}\{e^{-S_{\mathscr C}/k}g_j:1\le j,k\le M\}\)
form a form core.  On every fixed finite-dimensional \(E\subset D(\EK)\)
the omitted Euler tail satisfies
\(\|N_E^{-1/2}(\mathcal T_X)_EN_E^{-1/2}\|\to0\).  A finite omitted block
containing two distinct primes has positive definite Gram form on every
finite centered space.
\end{theorem}

The last assertion uses only that \(\log p/\log r\notin\Q\), so the group
generated by the two displacements is dense; a vector killed by both is
invariant under a dense subgroup, hence constant, hence zero after
centering.  Theorem~\ref{thm:heat} removes weighted spectral synthesis of
the zero modes as a separate hypothesis.

\begin{theorem}[Bounded prime generator]
\label{thm:bounded}
The literal ordinary-prime jump generator is bounded on the Doob space.
Consequently every normalized subthreshold eigenstate of the complete
prime--Gamma operator lies in the Gamma operator domain and has a uniform
local second logarithmic Fourier moment.
\end{theorem}

\subsection{The Gamma absorption theorem}
\label{subsec:absorption}

The following is the sharpest unconditional gain currently available and
is used in Entries~C5 and G2 below.  Set
\[
 r_\Gamma(u)=\frac{e^{-5u/2}}{1-e^{-2u}}\quad(u>0),
 \qquad
 a=\frac Kh,
 \qquad
 \mathcal C=T_{K'+K/2}M_a ,
\tag{3.1}
\]
where \(\mathcal C\) is the physical Abel connection of the delay-bank
realization.  For a nonnegative even displacement density \(v\) write
\[
 \mathfrak b_v(q)=\tfrac12\iint_{\R^2}v(x-y)K(x)K(y)|q(x)-q(y)|^2\dd x\dd y .
\tag{3.2}
\]

\begin{theorem}[Pointwise theta--Gamma margin]
\label{thm:margin}
For all \(u>0\),
\[
 \boxed{2K(u)<\frac{1501}{2000}\,r_\Gamma(u)},
 \qquad\text{hence}\qquad
 w_\Gamma:=r_\Gamma-2K\ \ge\ \frac{499}{2000}r_\Gamma>0 .
\tag{3.3}
\]
\end{theorem}

\begin{proof}
Put \(t=e^{2u}\ge1\).  From (2.2) and (3.1),
\[
 \frac{2K(u)}{r_\Gamma(u)}
 =4\pi t^{3/2}(t-1)\sum_{m\ge1}m^2(2\pi m^2t-3)e^{-\pi m^2t}.
\tag{3.4}
\]
Write \(R_1(t)=4\pi t^{3/2}(t-1)(2\pi t-3)e^{-\pi t}\) for the first
atom.  Its logarithmic derivative
\(d(t)=\frac3{2t}+\frac1{t-1}+\frac{2\pi}{2\pi t-3}-\pi\) satisfies
\(d'(t)<0\), so \(R_1\) has a single global maximum.  Using
\(\frac{333}{106}<\pi<\frac{355}{113}\) and monotonicity of \(d\) in
\(\pi\), rational arithmetic gives \(d(169/100)>0>d(17/10)\), so the
maximizer lies in \((169/100,17/10)\).  On that interval
\(t^{3/2}<\frac{17}{10}\cdot\frac{131}{100}\) since
\((131/100)^2>17/10\); also \(t-1<7/10\) and
\(e^{-\pi t}<e^{-56277/10600}<1/202\).  Hence
\[
 R_1(t)<4\cdot\tfrac{355}{113}\cdot\tfrac{17}{10}\cdot\tfrac{131}{100}
 \cdot\tfrac7{10}
 \Bigl(2\cdot\tfrac{355}{113}\cdot\tfrac{17}{10}-3\Bigr)\tfrac1{202}
 =\tfrac{240179723}{322417250}<\tfrac34 .
\tag{3.5}
\]
For \(m\ge2\), discard \(-3\) and use
\(\frac{(m+1)^4e^{-\pi(m+1)^2t}}{m^4e^{-\pi m^2t}}\le(\tfrac32)^4e^{-5\pi}<10^{-5}\)
to get \(\sum_{m\ge2}m^4e^{-\pi m^2t}<17e^{-4\pi t}\).  With \(s=t-1\) and
\((1+s)^{5/2}\le e^{5s/2}\),
\[
 t^{5/2}(t-1)e^{-4\pi t}\le\frac{e^{-4\pi}}{e(4\pi-5/2)} .
\tag{3.6}
\]
Using only \(\pi^2<10\), \(\pi>3\), \(e>2\) and \(e^{4\pi}>e^{12}>160000\),
the tail contributes less than \(1360/(19\cdot160000)<1/2000\).  Adding
(3.5) proves (3.3).
\end{proof}

\begin{remark}
The proof uses no floating-point input.  The true maximum of (3.4) is
\(\approx0.71779\), attained at \(t\approx1.69610\), so the stated
constant has genuine slack; the tightest step is
\(e^{-\pi\cdot169/100}<1/202\), where the true value is \(1/202.24\).
Any application needing only a constant below \(1\) may replace \(3/4\)
by \(4/5\) throughout and avoid that step.
\end{remark}

\begin{theorem}[Exact Gamma--connection factorization]
\label{thm:absorption}
Let \(q\) lie in the common form core and put \(F=hq\), \(f=Kq\).  Then
\[
 \boxed{
 \mathfrak b_{\Gamma,*}(q)+2\operatorname{Re}\langle F,\mathcal CF\rangle_{\omega_K}
 =\mathfrak b_{w_\Gamma}(q)+2\int_\R(K*K)(x)K(x)|q(x)|^2\dd x .}
\tag{3.7}
\]
Both terms on the right are nonnegative, and consequently
\[
 \mathfrak b_{\Gamma,*}(q)+2\operatorname{Re}\langle F,\mathcal CF\rangle_{\omega_K}
 \ \ge\ \frac{499}{2000}\,\mathfrak b_{\Gamma,*}(q).
\tag{3.8}
\]
\end{theorem}

\begin{proof}
Since \(r_\Gamma=w_\Gamma+2K\) we have
\(\mathfrak b_{\Gamma,*}=\mathfrak b_{w_\Gamma}+2\mathfrak b_K\).  The
adjoint computation for \(\mathcal C\) together with \(c_K=1/2\) gives
\(\operatorname{Re}\langle F,\mathcal CF\rangle_{\omega_K}=\langle f,T_Kf\rangle_2\).
Expanding the difference square in \(\mathfrak b_K\) and using
\(\int K(x-y)K(y)\dd y=(K*K)(x)\) yields the exact identity
\[
 \mathfrak b_K(q)+\langle f,T_Kf\rangle_2
 =\int_\R(K*K)(x)K(x)|q(x)|^2\dd x .
\tag{3.9}
\]
Combining the three displays proves (3.7); Theorem~\ref{thm:margin} and
monotonicity of \(v\mapsto\mathfrak b_v\) give (3.8).
\end{proof}

The argument is performed entirely at the difference-form level and never
splits the singular Gamma kernel into separately divergent diagonal and
convolution pieces.  Its consequence is that the Hermitian part of the
connection is no longer an open sign: it is absorbed by the completed
Gamma remainder with the explicit margin \(499/2000\), on every heat row,
every hybrid row, and indeed the full common core.  The \(K'\) part of
\(\mathcal C\) is exactly skew-adjoint and therefore lossless.

\subsection{The common PNT Stieltjes form}

Let
\[
 \mathcal D(u)=\sum_{\log n\le u}\frac{\Lambda(n)}{\sqrt n}-2(e^{u/2}-1),
 \qquad
 k_D(u)=e^{-u/2}\{\psi(e^u)-e^u+1\} .
\tag{3.10}
\]

\begin{proposition}[Common primitive]
\label{prop:pnt}
As distributions on \((0,\infty)\),
\(\dd\mathcal D=\dd k_D+\frac12k_D(u)\dd u\).  Consequently the outgoing
prime-number-theorem quadrature and the incoming Abel flux are the two
halves of a single Stieltjes form,
\[
 \boxed{
 \mathfrak P_{\rm PNT}(q)
 =\int_{[0,\infty)}\mathcal J_u(q)\dd\mathcal D(u)
 =-\int_0^\infty k_D(u)\{\mathcal J_u'(q)-\tfrac12\mathcal J_u(q)\}\dd u .}
\tag{3.11}
\]
\end{proposition}

Together with Theorem~\ref{thm:absorption}, the exact remaining physical
form is
\[
 \mathfrak P_{\rm PNT}(q)+\mathfrak b_{w_\Gamma}(q)
 +2\int_\R(K*K)K|q|^2\dd x\ \ge\ 0,
\tag{3.12}
\]
with the simpler sufficient version
\[
 \boxed{\mathfrak P_{\rm PNT}(q)+\tfrac{499}{2000}\mathfrak b_{\Gamma,*}(q)\ \ge\ 0 .}
\tag{3.13}
\]
Entry~C5 shows that (3.12)--(3.13) cannot be proved by any
displacement-local mechanism.

\section{The circularity boundary}
\label{sec:circularity}

This section states the organizing theorem of the paper.  It explains why
the catalogue is the correct output, and it is also the sharpest single
constraint on any future attempt.

\begin{theorem}[Off-line anti-surplus]
\label{thm:antisurplus}
Assume \(\XiR\) has an off-line zero orbit.  Then there is a real even
\(F\in\NK\) in the common form domain, \(q=F/h\), such that for every
finite prime-power head \(X\):
\begin{enumerate}
\item the finite-head pivot is negative,
 \(\mathcal A_X(q,q)=-\delta_X<0\);
\item the exact adaptive gain is the complete omitted-prime response,
 \(G_X=\mathcal T_X(q,q)\);
\item the source surplus is strictly negative,
\[
 \boxed{G_X-\delta_X=\mathcal A_\infty(q,q)<0 ;}
\tag{4.1}
\]
\item for every finite tail cutoff \(Y>X\) the physical bordered
 determinant is negative,
 \(\mathfrak C_X(Y)-\delta_X<0\).
\end{enumerate}
All displacement forms here are the literal maps
\(q\mapsto\sqrt{K(x)K(x-\log n)}\{q(x)-q(x-\log n)\}\) with the ordinary
weights \(\Lambda(n)/\sqrt n\), with the complete Gamma channel, the
polar threshold, and the equation \((hq)*K=0\) all retained jointly.
\end{theorem}

\begin{proof}
An off-line orbit produces a real even \(F\in\NK\) with
\(\mathcal A_\infty(q,q)<0\).  The exact decomposition
\(\mathcal A_\infty=\mathcal A_X+\mathcal T_X\) and nonnegativity of every
omitted atom give (1).  In the scalar row there is no old-mode regression
and no radical shorting, so the resolvent defining the adaptive gain is
the identity, giving (2); adding the two identities gives (3).
Nonnegativity of the omitted atoms gives \(\mathfrak C_X(Y)\le G_X\), and
with no nuisance coordinate the deleted principal determinant is \(1\),
which gives (4).
\end{proof}

\begin{corollary}[Counterfactual invariance]
\label{cor:invariance}
The positive data
\[
 \Lambda(p^k)>0,\quad K>0,\quad \mathcal J_{\log p^k}(q,q)\ge0,
 \quad (hq)*K=0
\]
all hold for the vector constructed under the off-line counterfactual,
while the strict surplus \(G_X>\delta_X\) does not.  Hence no implication
which is insensitive to the location of the mean-periodic frequencies can
prove the surplus.
\end{corollary}

\begin{remark}[The surplus is not a lemma]
\label{rem:notlemma}
Identity (4.1) says that the desired surplus is not a preliminary
estimate from which completed positivity follows.  It is completed
positivity, evaluated at the stationary residual.  Any proposed proof
should therefore be audited at the first line where (4.1) is assigned a
positive sign; that line carries the entire Riemann-strength input.  This
is the single most useful practical consequence of the present work.
\end{remark}

\section{Taxonomy}

Each entry below carries one of four tags.

\begin{description}
\item[\NOGO] A theorem in the literal Riemann system showing that a class
 of arguments cannot produce (2.7).  Strongest form.
\item[\FALS] A construction, valid under the off-line counterfactual and
 inside the literal system, on which the proposed inequality is false.
 By Corollary~\ref{cor:invariance} such an entry shows the mechanism
 cannot be counterfactual-insensitive.
\item[\CTMD] An explicit model satisfying every structural hypothesis of a
 proposed mechanism and violating its conclusion.  Rules out proofs using
 only the listed hypotheses; does not itself concern \(\zeta\).
\item[\SCAL] A quantitative envelope mismatch: the mechanism is valid but
 provably loses more than the quantity it must control.
\end{description}

We emphasize the difference between \NOGO{} and \CTMD.  A countermodel
rules out an \emph{axiomatization}: it shows that the properties invoked
are insufficient.  It leaves open that the actual Riemann system has
additional properties which close the argument.  A no-go rules out the
mechanism in the actual system.

\section{Family A. Order, positivity preservation, total positivity}

\begin{entry}[\NOGO{} Failure of the first Beurling--Deny criterion]
\label{ent:A1}
In the completed source formula, let \(\varphi_1,\varphi_2\ge0\) be
disjointly supported bumps whose separation \(d\) avoids every prime-power
logarithm.  Their cross form is strictly positive: the polar kernel grows
like \(e^{d/2}\) while the opposing Gamma kernel decays like \(e^{-d/2}\).
Hence the completed heat semigroup is not positivity preserving, in the
full sector and in the even sector.
\end{entry}

\emph{Consequence.} No Perron--Frobenius, Harnack, or ground-state
domination argument based on the natural pointwise cone applies to the
completed operator.  This is significant because the ground-state
heuristic --- ``the bottom eigenvector should be positive, hence the
radical should be the bottom'' --- is the most natural first attack, and
it is unavailable.

\begin{entry}[\NOGO{} The system is not \(\mathrm{TP}_2\)]
\label{ent:A2}
The exact Gamma density \(g\) is strictly log-convex, so
\(g(t)^2-g(t-\eps)g(t+\eps)<0\).  Hence neither the Gamma kernel nor the
short-time full semigroup is totally positive of order two.
\end{entry}

\emph{Consequence.} Variation-diminishing spectral ordering is
unavailable.  In particular one cannot conclude from the fact that the
threshold state \(K''/K\) has a single sign change that it is the first
centered mode.

\begin{entry}[\NOGO{} Failure of stochastic monotonicity]
\label{ent:A3}
The folded ordinary-prime--Gamma Doob process has, for each displacement
\(u\), the reflected branch \(t\mapsto|t-u|\).  With the literal atom
\(p=2\), at \(t_0=\frac14\log2\) and \(z=\frac34\log2\), the upper-tail
jump rate has a downward discontinuity of exactly
\(c_K\Lambda(2)K(z)/(\sqrt2\,h(t_0))>0\); the Gamma continuum and every
other prime-power atom form a continuous remainder because no other
displacement meets that boundary.  This violates the necessary tail-rate
order.
\end{entry}

\emph{Consequence.} Both the short-time semigroup and the
large-parameter resolvent reverse the increasing cone on explicit smooth
generator tests.  Ordering \(K''/K\) as the first centered mode by
stochastic monotonicity or by an invariant increasing cone is impossible.
Together with A2, the entire ``order'' family is closed.

\section{Family B. Coupling, transport, martingale routing}

\begin{entry}[\NOGO{} No radical-sharp first-chaos coupling]
\label{ent:B1}
Restore the continuous theta coordinate and the integer and rational
theta marks in all four channels (Gamma, divisible, nondivisible
fractional, central crossing), giving an exact positive latent Dirichlet
form with the literal coefficients.  No positive conditional-expectation
or first-chaos projection maps this lift to the polar pair gradient while
preserving the full equality family
\(\{K^{(2j)}/K\}\).
\end{entry}

\emph{Mechanism.} Equality in conditional Jensen makes every radical edge
signature deterministic; the signatures separate oriented endpoints and
therefore force endpoint preservation.  But the literal prime edges carry
singular source mass on the lines \(t-s=\log n\) and \(t+s=\log n\), while
the polar pair law is absolutely continuous.  This is an exact
measure-capacity contradiction.  A second, independent capacity
obstruction holds in the Gamma channel alone: its density lacks polar
capacity on remote fixed-displacement rectangles.

\begin{entry}[\NOGO{} No positive martingale routing]
\label{ent:B2}
In the folded polar theta law with dilation variable \(S=Be^T\), choosing
a prime-power divisor \(N\mid b\) with probability \(\Lambda(N)/\log b\)
identifies the divisor current and the cross-divisor dispersion as the
conditional mean and variance of one increment, and the resulting ANOVA
transform is unitary.  Nevertheless every finite or countable positive
martingale routing retaining all currents, dispersions, fractional and
central fibres, and Gamma, obeys an exact nonnegative defect identity;
radical equality forces every martingale difference and every unused
source square to vanish, hence endpoint preservation, hence the
contradiction of B1.  The natural separated assignment already fails on
the two-state fibre \(2<S<3\).
\end{entry}

\begin{entry}[\NOGO{} Shorting is not a weakening]
\label{ent:B3}
Polarization shows that shorting subtracts exactly the same radical square
from the source and from the target.  Hence the complementary inequality
is \emph{equivalent} to, not weaker than, the original all-prime
inequality.  Moreover radical completeness is refuted exactly: for every
zero \(z\), the mode \(\chi_z(x)=\cos(zx)/\cosh(x/2)\) is a nonzero
\(L^2(\mu_K)\) element orthogonal to all \(K^{(2j)}/K\).
\end{entry}

\emph{Consequence.} One cannot buy room by passing to the radical
complement.  Every reformulation obtained by shorting inherits the exact
original difficulty.  This entry is used repeatedly below.

\section{Family C. Local squares, sums of squares, Hodge--Korn}

This family contains the master obstruction of the paper.

\begin{entry}[\NOGO{} Radical saturation: no local positive remainder]
\label{ent:C1}
The family \(r_j=K^{(2j)}/K\), \(j\ge1\), satisfies
\[
 \EK(r_j)=\tfrac12\Var_{\mu_K}(r_j)\qquad(j\ge1),
\tag{7.1}
\]
so the sharp inequality (2.7) is saturated on an infinite-dimensional
space.  Suppose a proposed proof has the form
\[
 \EK(r)-\tfrac12\Var_{\mu_K}(r)=\sum_\alpha\mathcal S_\alpha(r),
 \qquad \mathcal S_\alpha\ge0,
\tag{7.2}
\]
and that one summand dominates a local edge square,
\[
 \mathcal S_{\alpha_0}(r)\ \ge\
 \int_0^\infty\!\!\int_\R w(x,u)|r(x)-r(x-u)|^2\dd x\dd u
\tag{7.3}
\]
with \(w\ge0\) positive on a set of positive measure.  Then (7.2) is
impossible.
\end{entry}

\begin{proof}
Take \(r_1=K''/K\).  By (7.1) the left side of (7.2) vanishes, so every
nonnegative summand vanishes; (7.3) then forces \(r_1(x)=r_1(x-u)\) for
almost every \((x,u)\) in the positive set of \(w\).  As \(r_1\) is real
analytic, Fubini and analytic continuation make it periodic for at least
one nonzero displacement.  But the leading positive theta term gives
\(K''(x)/K(x)\asymp e^{4x}\) as \(x\to+\infty\), so \(r_1\) is unbounded
and cannot be periodic.
\end{proof}

\emph{Consequence.} The following are all excluded:
\begin{enumerate}
\item canonical-path comparison using only divisible theta indices;
\item Efron--Stein decomposition over individual prime towers;
\item any Cheeger estimate with a strict edge surplus;
\item any sum-of-squares proof whose individual squares are local
 translation differences.
\end{enumerate}
What survives is a \emph{signed} global decomposition with genuine
cancellation between different prime channels.  Entry~C1 is the reason
every later entry in this paper insists on the word ``signed''.

\begin{entry}[\NOGO{} Finite-incidence Hodge--Korn rigidity]
\label{ent:C2}
The vertex signatures \((K^{(2j)}(x)/K(x))_{j\ge0}\) are finitely linearly
independent.  Hence a local linear combination of edge increments which
vanishes on the complete radical is divergence free; it is an incidence
cycle and therefore vanishes on every gradient.  Since the theta-divisor
fibres are stars, every radical-sharp local amplitude on their current,
dispersion or residue-character coordinates is zero.
\end{entry}

\emph{Consequence.} No nontrivial integral of finite-local squares can
complete the signed Gamma--polar channel.  After the exact global radical
anti-short one obtains the sharp identity
\(\QW=\|G_\perp r\|^2-\|D_\perp r\|^2\), whose remaining Korn contraction
is globally nonlocal and is violated by a hypothetical subthreshold mode
by the factor \((2\alpha)^{-1/2}>1\).

\begin{entry}[\NOGO{} No infinite boundary escape]
\label{ent:C3}
In the literal source norm, with every von Mangoldt and theta weight
retained, the conductance of the folded spatial cut at \(R\) is
\(O(\exp\{-ce^R\})\).  Compact smooth even multipliers form a core of the
closed prime--Gamma form, so every bounded locally divergence-free current
is orthogonal to all form-domain gradients and has zero flux at the folded
end.  A nonzero cofinal boundary flux would require squared source norm at
least \(\exp\{ce^R\}\).
\end{entry}

\emph{Consequence.} C2 cannot be evaded by moving the missing amplitude to
infinity.  There is no extra bounded boundary-homology sector.

\begin{entry}[\NOGO{} Primitive and Volterra squares are indefinite]
\label{ent:C4}
Pushing the identities \(k_y=\frac12(D^2-\frac14)f_y\) and
\(\phi(y)=-\pi(y^3e^{-\pi y^2})'\) through the complete source, the
multiplicative pushforward is the single literal signed measure
\[
 \dd\eta=\sum_{n\ge2}\Lambda(n)\delta_n
 +\Bigl\{-1+\frac1{z(z^2-1)}\Bigr\}\dd z .
\tag{7.4}
\]
Fibrewise positivity is false: before the first prime, \(\eta\) is
strictly negative on \((4/3,2)\), and even the quadratic increment
satisfies \(A_{(\log z)^2}(2^-)<-\frac{43}{4500}\).  Summing every theta
and character fibre gives exactly \(F_r(z)=z^{-1/2}J_{\log z}(r)\), whose
cumulative rises to the plastic root, then has strictly negative
continuous drift between prime powers with upward jumps
\(\Lambda(n)n^{-1/2}J_{\log n}\).
\end{entry}

\emph{Consequence.} Neither the fibrewise nor the fully aggregated
coordinate is a local positive square.  The endpoint terms vanish, so
Gamma and the pole cannot be appended afterwards as a boundary
correction: they are already inside the signed cumulative.

\begin{entry}[\NOGO{} No centered-path or local Stieltjes completion]
\label{ent:C5}
Let \(C_uq=\{K(x)K(x-u)\}^{1/2}\{q(x)-q(x-u)\}\) and
\(Z_q(u)=e^{-u/4}C_uq\), so \(\|Z_q(u)\|^2=e^{-u/2}\mathcal J_u(q)\).  For
\(0<\theta\le1\) set \(a_\theta(u)=\theta e^{u/2}r_\Gamma(u)>0\).  Then
\[
 \mathfrak P_{\rm PNT}+\theta\mathfrak b_{\Gamma,*}
 =\int_0^\infty a_\theta\Bigl\|Z_q-\frac{E_1(e^u)}{a_\theta}Z_q'\Bigr\|^2\dd u
 -\int_0^\infty\frac{E_1(e^u)^2}{a_\theta(u)}\|Z_q'(u)\|^2\dd u ,
\tag{7.5}
\]
with \(E_1(T)=\psi(T)-T+1\).  The second line has the wrong sign and is
not an omitted Gamma term: the entire Gamma budget is already inside
\(a_\theta\).  Moreover, at the available margin \(\theta_*=499/2000\),
the scalar density of the sufficient form (3.13) is
\(w_*(u)=-e^{u/2}+\theta_*r_\Gamma(u)\), and with \(z=e^u\)
\[
 \frac{\theta_*r_\Gamma(u)}{e^{u/2}}=\frac{\theta_*}{z^3-z},
 \qquad
 z\ge\tfrac98\ \Longrightarrow\ z^3-z\ge\tfrac{153}{512}>\tfrac{499}{2000},
\tag{7.6}
\]
since \(306000>255488\).  Hence
\[
 \boxed{w_*(u)<0\qquad\bigl(\log\tfrac98\le u<\log2\bigr).}
\tag{7.7}
\]
In the exact residual (3.12) the density \(-e^{u/2}+r_\Gamma-2K\) is
likewise strictly negative on \([\log\frac32,\log2)\), because
\(z^3-z\ge\frac{15}8>1\) and \(K>0\).
\end{entry}

\emph{Quantitative version.}  Let \(q_N=\chi\cos(Nx)\) with
\(0\ne\chi\in C_c^\infty\) chosen so that its support overlaps its
translates on a fixed \(I\Subset(0,\log2)\).  Then
\(\int_I\|Z_{q_N}'(u)\|^2\dd u\ge c_IN^2\) for large \(N\), whereas
\[
 \mathfrak b_{\Gamma,*}(q_N)=O(\log N),
 \qquad
 \mathfrak P_{\rm PNT}(q_N)=O(1).
\tag{7.8}
\]
The two \(N^2\) terms in (7.5) cancel each other exactly; estimating
either separately destroys the cancellation by a factor larger than any
fixed Gamma margin.  Projection off any fixed finite radical block changes
these quantities by \(O(N^{-A})\) for every \(A\).

\emph{Consequence.} None of the following can be the next step: positivity
of the common displacement density; a pointwise completion of the centered
Abel derivative; a bound of \(\mathcal J_u'\) or \(Z_q'\) by the Gamma
form; or a finite-radical version of any of these.  Note that (7.7)
locates the failure \emph{below the first prime}, where the arithmetic is
entirely explicit --- the obstruction is not an artifact of prime
irregularity.

\section{Family D. Determinants, minors, charging}

\begin{entry}[\NOGO{} Cauchy--Binet charging is circular]
\label{ent:D1}
In exterior-power coordinates, an abstract contractive charge exists
exactly when \(\tau_{d+1}\ge\delta_J\tau_d\), i.e. exactly when the
inequality to be proved holds.  Consequently the squared-minor expansion
supplies no sign.  Moreover a universal positive endpoint-local
weight-preserving charge would saturate Jensen on every radical signature,
hence preserve endpoints, and then contradict B1.
\end{entry}

\emph{Consequence.} The only surviving possibility is a row-dependent,
globally nonlocal charge after the complete radical anti-short, whose
strict norm bound is precisely the projected Pl\"ucker-frame inequality.
That is the surplus again, not a route to it.

\begin{entry}[\CTMD{} Network positivity does not force crossing]
\label{ent:D2}
There is an exact four-state rational model with a connected full-rank
tail, nontrivial adaptation, an exact threshold radical, positive Gamma
and head graph channels, literal weights \(\Lambda(n)/\sqrt n\), and a
theta-sized envelope, for which the bordered determinant remains negative.
\end{entry}

\emph{Consequence.} Effective resistance, matrix-tree expansions, generic
network positivity and summability do not force the crossing.

\begin{entry}[\CTMD{} Exact saturation with the literal theta data]
\label{ent:D3}
A one-dimensional model using the actual theta displacement forms, the
actual lengths \(\log p^k\), and the ordinary weights
\(\Lambda(p^k)=\log p\), with the free source diagonal set equal to the
complete positive tail, has every proper determinant strictly below the
deficit, converging to it.  Permanent non-crossing is exactly
\(\sigma_\infty\le0\).
\end{entry}

\emph{Consequence.} This rules out any proof from source algebra, positive
ordinary-prime increments, inertia and summability alone.  Total
positivity supplies no missing input either: the Gamma kernel is not
\(\mathrm{TP}_2\) (A2), the literal prime tower has a \(\mathrm{PF}_2\)
failure, and the polarized jet \(j_2\) is already indefinite on
\(\{p,p^2\}\).

\begin{entry}[\SCAL{} The observation-only determinant is strictly lossy]
\label{ent:D4}
The pure Gram ratio
\(\det G_{\mathcal B}^{(M)}/\det G_{\mathcal B}^{(M-1)}\) is a valid but
strictly weaker certificate than the exact adaptive gain.  In stable
floating-point diagnostics at dimensions \(4,7,12\) it misses the initial
deficit while the exact Kalman gain crosses it, sometimes by five orders
of magnitude.  The missing energy is the component parallel to the old
observation range, priced by the positive preceding prime--Gamma block.
\end{entry}

\emph{Consequence.} The force-bearing object is the augmented determinant
\[
 \Delta_{\mathcal B}
 =\frac{\det\begin{pmatrix}A+U^*U&U^*r\\ r^*U&\|r\|^2\end{pmatrix}}
        {\det(A+U^*U)},
\tag{8.1}
\]
not the observation Gram ratio.  We record this because the weaker
certificate is the one that suggests itself first.

\section{Family E. Finite heads, staircases, adaptive filtering}

\begin{entry}[\NOGO{} No finite-head telescope]
\label{ent:E1}
Every proper finite prime-power head is strictly negative on each
nonconstant exact Riemann-radical direction.  Hence its signed short over
the radical is unbounded below, and the full infinite prime--Gamma--polar
limit must be assembled \emph{before} radical shorting.
\end{entry}

\emph{Consequence.} A finite-head Cholesky or Pick telescope is
impossible.  This is the finite-level shadow of B3.

\begin{entry}[\CTMD{} Discrete levels and full-rank sensors are not enough]
\label{ent:E2}
There is an exact three-state reversible graph generator with a constant
ground state, an exact completed \(1/2\)-radical, discrete levels, and an
infinite norm-convergent bank of strictly positive full-rank sensors
carrying the literal weights \(\Lambda(p^k)/p^{k/2}\), which nevertheless
retains a complementary eigenvalue \(1/4\).
\end{entry}

\begin{entry}[\CTMD{} Adaptation loss defeats tail energy]
\label{ent:E3}
There is an injective two-mode theta-scale model whose residual tail energy
strictly exceeds its deficit, but which loses enough energy to adaptive
regression that its completed pivot remains negative.
\end{entry}

\emph{Consequence.} Schur adaptivity, positive literal weights, strict
observability and summability do not force crossing.  The missing input is
an estimate of the adaptation-loss term itself.

\begin{entry}[\CTMD{} Exact equality saturation over the prime powers]
\label{ent:E4}
There is a saturation model indexed by the ordinary prime powers, with the
literal weights \(\Lambda(m)/\sqrt m\), an \(e^{-2\pi m}\) theta envelope,
injective atom features and zero adaptive leakage, whose cumulative gain
equals \(\delta\) exactly: every finite pivot is negative and the
completed pivot is zero.
\end{entry}

\emph{Consequence.} Stieltjes positivity and theta summability perform the
finite selection only \emph{after} a strict surplus has been proved.
Equality does not suffice.  This entry is the reason the target inequality
must be strict, and the reason a compensated vanishing-tolerance form is
needed for the limit.

\begin{entry}[\NOGO{} Response--inertia bound]
\label{ent:E5}
For a Hermitian \(H\) and any \(T\),
\(\nu_-(H|_{\ker T})\ge\nu_-(H)-\operatorname{rank}T\).  Hence whenever a
signed finite spectral block has more negative directions than available
prime-midpoint rows, some negative vector annihilates every one of those
rows.
\end{entry}

\emph{Consequence.} No arbitrary-vector midpoint frame floor may be used.
For the actual staircase the dimension obstruction disappears, because
\(A\succ0\) with a negative Schur pivot forces negative index exactly one;
but no inertia identity then bounds the regression residual from below.

\begin{entry}[\NOGO{} Local observability does not pay the deficit]
\label{ent:E6}
On every finite analytic even mode space the literal displacement feature
restricted to any nonempty midpoint aperture is injective, giving the
computable Christoffel bound
\(\mathcal J_{u,W}(q^*)\ge\kappa_M(u,W)>0\).  However, for a nonzero
completed radical vector the finite-head deficit equals the sum of
\emph{all} omitted positive atom energies.  Hence every proper finite
omitted block gains strictly less than the deficit, and its
moving-midpoint mass can tend to zero.
\end{entry}

\emph{Consequence.} Any surviving estimate must use the completed
anti-short \(q^*\in\mathcal R^\perp\) quantitatively; local observability,
the adaptive source equation, and tail localization are individually
insufficient.

\begin{entry}[\SCAL{} Prime-sampling envelope mismatch]
\label{ent:E7}
On every fixed spectral block \(|\operatorname{Im}z|\le b<1/2\) with jets
of order at most \(J\), and every fixed \(\delta>0\),
\[
 \sum_{X<p\le(1+\delta)X}\beta_p|A_p(q)+\rho_p(q)|^2
 \ \ge\ c'_{\mathcal Z,\delta}X^{3-b}(1+\log X)^{-2J}
 e^{-2\pi(1+\delta)X}\|q\|_{\rm coeff}^2 .
\tag{9.1}
\]
This is a genuine quantitative sampling theorem, proved from the prime
number theorem, but compared with the natural tail scale \(e^{-2\pi X}\)
it loses \(e^{-2\pi\delta X}\).
\end{entry}

\emph{Consequence.} A relative-interval attack cannot place the response
inside the \(O(1)\) prime window where the theta weight is comparable to
its leading tail value.  Remez and Tur\'an localization on the same
relative interval do not repair this.  Any successful estimate must be a
short-window statement.

\section{Family F. Spectral and operator-theoretic detectors}

\begin{entry}[\NOGO{} Birman--Schwinger supplies no sign]
\label{ent:F1}
A compactly supported positive boost produces a compact
Birman--Schwinger operator \(\mathcal K_{\lambda,R,M}\) whose eigenvalue
\(1\) detects each isolated subthreshold bound state.  But the exact
Rayleigh quotient shows that \(\|\mathcal K_{\lambda,R,M}\|\le1\) is,
term for term, equivalent to \(\EK\ge\lambda\|\cdot\|^2\).
\end{entry}

\emph{Consequence.} Localization proves Fredholm compactness and supplies
no sign.  This is the faithful nonlocal version of the earlier
Birman--Schwinger wall, not a new mechanism.

\begin{entry}[\NOGO{} Mourre and dilation commutators]
\label{ent:F2}
After ground-state conjugation, the translation commutator is parity odd
and has identically zero quadratic form on the physical even sector.
Dilation is parity even, but every prime power at \(a=\log n\) contributes
the isolated order-one distribution \(a(\delta_a'-\delta_{-a}')\), and the
complete raw dilation commutator is unbounded in both signs on compact
oscillatory even rows.  In the exact complement, translation is tangent to
\(F*K=0\) but changes parity, while dilation preserves parity and has
transverse leakage \(t\XiR'(t)\widehat F(t)\) at every simple zero.  Below
the essential threshold a strict projected Mourre estimate holds exactly
when the corresponding spectral projection is zero, and the normalized
commutator vanishes on heat or hybrid rows concentrating on a bound
cluster.
\end{entry}

\begin{entry}[\NOGO{} Heat moments and operator-monotone transforms]
\label{ent:F3}
Every raw heat Hankel matrix is automatically positive, whereas eventual
positivity of any one odd moment shifted to the threshold is already
equivalent to the physical surplus.  The same wall persists after every
strictly increasing operator-monotone transform.
\end{entry}

\emph{Consequence.} Heat summability and scalar moment positivity supply
no hidden fraction of the missing bound.  Relatedly, the relative-heat
observable has logarithmic decay rate exactly
\(-\{\inf\sigma(L|_{\mathscr C})+\frac12\}\), so the desired surplus is
equivalent to rate at most \(-1\) while the unconditional energy gives
only \(-1/2\).

\begin{entry}[\NOGO{} Positive source atoms tolerate a subthreshold mode]
\label{ent:F4}
If a subthreshold ground cluster \(\alpha<1/2\) of multiplicity
\(m_\alpha\) exists, the Birman--Solomyak relative measure has exact local
density \(m_\alpha\dd\lambda\), so
\(\Theta_V(t)\sim m_\alpha e^{-t(\alpha+1/2)}\) for every positive
injective trace-class boost, and the heat-weighted source satisfies
\(\frac12A_V-E_V\sim(\frac12-\alpha)m_\alpha t^{-1}e^{-t(\alpha+1/2)}>0\).
\end{entry}

\emph{Consequence.} Positive source atoms and all unshifted heat moments
remain compatible with a subthreshold mode.  A closure must estimate the
\emph{compensated} source itself.

\begin{entry}[\SCAL{} Non-uniformity as the deficit vanishes]
\label{ent:F5}
For every fixed \(\delta>0\) the subthreshold spectral projection below
\(1/2-\delta\) is finite rank, and one finite heat-core compression detects
every such eigenstate below \(1/2-\delta/2\).  The reduction is not
uniform as \(\delta\downarrow0\): packets of amplitude \((\log N)^{-1/2}\)
have vanishing Hilbert mass at frequency \(N\) but fixed Gamma form
energy, and Hilbert accuracy \(\eta^2\asymp\delta\) already requires
\(\Omega\gtrsim\exp(C_R/\delta)\).
\end{entry}

\emph{Consequence.} Fixed-gap violations genuinely reduce to finite
arithmetic; near-threshold violations do not.  Any finite-dimensional
computational strategy must therefore come with a proof that the deficit
is bounded away from zero, which is itself part of the problem.

\begin{entry}[\NOGO{} No central tightness from mean periodicity]
\label{ent:F6}
Multiplicity jets are essential to local \(\XiR\)-universality: a
compactly supported annihilator would have an entire Fourier--Laplace
transform with \(O(T)\) zeros, while Riemann--von Mangoldt gives
\(\asymp T\log T\) zeros with multiplicity.  Hence exact mean-periodic
zero modes approximate zero on any central window and an arbitrary profile
on a disjoint window.  A diagonal construction gives normalized real-even
graph-domain vectors in the exact radical complement whose mass escapes
equally to both tails while every localized second-log norm tends to zero.
\end{entry}

\emph{Consequence.} Mean periodicity, evenness, and the regularity of
Theorem~\ref{thm:bounded} cannot yield central tightness.  The remaining
input must be a global threshold-resolvent or eigenstate-observability
estimate using the literal operator equation.

\section{Family G. Realization and factorization rigidity}

This family is the most structurally informative: it shows that the
missing constant is invariant under every repackaging of the problem.

\begin{entry}[\NOGO{} Canonical signed transfer has a universal threshold]
\label{ent:G1}
After constants and the complete radical have been shorted, the unique
exact signed transfer is
\(C_0=2^{-1/2}U_DA^{-1/2}U_A^*\), with
\[
 \boxed{\|C_0\|=(2\inf\sigma(A))^{-1/2}.}
\tag{10.1}
\]
Hence every exact signed kernel has the \emph{same} contraction threshold,
and a subthreshold eigenmode is amplified by precisely
\((2\alpha)^{-1/2}>1\).  A virial identity further shows that a square
plus a trace-zero commutator cannot conceal such a mode.
\end{entry}

\begin{entry}[\NOGO{} Compact all-atom factorization is unique-minimal]
\label{ent:G2}
Let \(\widetilde{\mathcal G}\) be the complete literal von
Mangoldt--Gamma displacement gradient and
\(\widetilde A=\widetilde{\mathcal G}^*\widetilde{\mathcal G}\) after exact
anti-shorting.  Both \(\mathcal J=[T_K,M_a]\) and
\(\mathcal C=T_{K'+K/2}M_a\) are Hilbert--Schmidt on the physical
mean-periodic space, with
\(\|\mathcal T_b\|_{\rm HS}^2=c_K^2\iint|b(x-y)|^2\dd\omega_K\dd\omega_K\).
Moreover \(\widetilde A\succeq\alpha_0I\) for some \(\alpha_0>0\), by
Theorem~\ref{thm:ess} together with triviality of the kernel after
removing constants.  Hence \(\mathcal C=H_C\widetilde{\mathcal G}\) with
\(H_C=\mathcal C\widetilde A^{-1/2}U_A^*\) compact and of minimal norm,
\[
 \inf_{\mathcal C=H\widetilde{\mathcal G}}\|H\|^2
 =\|\mathcal C\widetilde A^{-1}\mathcal C^\sharp\|,
 \qquad
 \boxed{\|H_C\|\le1\iff\mathcal C^\sharp\mathcal C\preceq\widetilde A .}
\tag{10.2}
\]
These gains are operator-norm limits of common-cutoff regularized
factorizations containing every prime power cofinally.
\end{entry}

\emph{Consequence.} A change of realization cannot manufacture the strict
gain.  Entries G1 and G2 together constitute a rigidity theorem about the
entire programme: the remaining constant is not a design parameter, and
any future reformulation will land on the same threshold.  We regard this
as the single most decision-relevant fact for anyone continuing the work.

\begin{entry}[\NOGO{} Local passivity cells cannot be multiplied]
\label{ent:G3}
The explicit positive periodic kernel
\(K_L=1+\eps\cos(4\pi x/L)+\eps\cos(6\pi x/L)\) has exact mean-periodic
row \(F_L=\cos(2\pi x/L)\), fixed Hermitian connection power
\(\eps^3/16\), and Gamma plus any fixed finite bank of literal delays
\((\log n,\Lambda(n)/\sqrt n)\) of size \(O(L^{-2})\).  The full minimal
Schur cost also stays bounded away from zero, and the witness is exactly
representable on a heat or hybrid Galerkin row.  A moving bank can avoid
the obstruction only if its weighted second delay moment is
\(\gtrsim L^2\), whereas an individual phase-resolving delay has
\(\log n\asymp L\).
\end{entry}

\emph{Consequence.} Derivative storage cannot create the missing sign, and
independently passive prime cells cannot be multiplied to obtain it.  This
is explicitly an abstract positive-kernel falsifier, not the theta kernel;
its content is that any certificate must use the cofinal prime
\emph{placement} jointly with Gamma, theta, and the polar channel, rather
than Gamma or a fixed prime head.

\begin{entry}[\NOGO{} Nonlinear amplification is unavailable]
\label{ent:G4}
On the exact complement \((hq)*K=0\): translation, M\"obius and Jordan
multipliers preserve each frequency and cannot amplify its transverse
exponent.  Pointwise powers do amplify it, but an off-axis mode
\(\cos(zx)^k\) necessarily leaves the zero divisor once
\(k|\operatorname{Im}z|>1/2\); tensor powers retain the equations only
before the non-intertwining diagonal restriction.  The spatial M\"obius
inverse cannot be passed through the zero-mode convolution: the required
squarefree absolute series diverges, and the regularized symbol is
\(\zeta(s)/\zeta(s+\eps)\), singular at exactly the hypothetical zero.
\end{entry}

\begin{entry}[\NOGO{} The bar homotopy has no critical contraction]
\label{ent:G5}
The normalized multiplicative bar homotopy
\(Z_\eps^{-1}=\sum_{k\ge0}(-X_\eps)^k\) recovers every literal von
Mangoldt weight from \(Z_\eps^{-1}\delta Z_\eps\).  If ordered
factorization chains remain orthogonal, the signed collapse on an
\(r\)-prime squarefree fibre has exact norm
\((\sum_kk!S(r,k))^{1/2}\ge\sqrt{r!}\); if they are cancelled before the
norm, the result is the M\"obius inverse itself.  On
\(q_\gamma=\cos(\gamma x)/\cosh(x/2)\), a critical zero of multiplicity
\(m\) forces homotopy gain \(1/|\zeta(\bar\rho+\eps)|\asymp\eps^{-m}\) and
connection gain \(m/\eps+O(1)\).  The specialized bar complex has nonzero
zeroth homology.
\end{entry}

\emph{Consequence.} Gamma and the pole are analytic and nonzero at the
zero, so completion preserves the residue.  Supplying the zero-fibre map
separately returns exactly to the canonical transfer of G1.

\section{Family H. The semilocal branch-selection coordinate}

The remaining entries concern the finite CCM coordinate rather than the
Doob coordinate, and are included because that is where the programme
began.

\begin{entry}[\NOGO{} Fixed derivative models miss off-line channels]
\label{ent:H1}
Every fixed span of cutoff derivatives of Riemann's full kernel is
asymptotically orthogonal to a hypothetical off-line evaluation channel,
since every transform in that span contains the common factor \(\XiR\).
In particular, coupling \(K\) and \(K'\) does not repair selection.
\end{entry}

\begin{entry}[\NOGO{} Many-vector co-Poisson frames do not improve the quotient]
\label{ent:H2}
If \(a_-(Q_L)\) is the lower frame bound of the low-mode
co-Poisson/prolate family on the negative spectral space, then
\[
 \sum_j\|P_-A_Lq_{j,L}\|^2\ \ge\ a_-(Q_L)
 \sum_{\epsilon_{\ell,L}<0}\epsilon_{\ell,L}^2 ,
\tag{11.1}
\]
so a fixed negative branch forces the aggregate overlap and residual to
vanish at the same rate.  The scale audit is also adverse: Gamma
coercivity and the absolute \(O(\lambda)\) prime--pole norm control at most
an \(O(Le^{C\lambda})\)-dimensional negative space, while the Slepian
family has only \(O(\lambda^2)\) modes.
\end{entry}

\begin{entry}[\NOGO{} Unquantified cyclicity diverges]
\label{ent:H3}
Although translates and derivatives of \(K\) are cyclic in ordinary
\(L^2\), an off-line zero \(\rho=\beta+i\gamma\) forces every truncated
radical approximant of a detecting test to have squared error
\(\gg\lambda^{-2(\beta-1/2)}\).  The semilocal lower-bound cost is
therefore \(\gg\lambda^{2(1-\beta)}\), which diverges rather than
vanishes.
\end{entry}

\begin{entry}[\NOGO{} A previously proposed lemma is false as stated]
\label{ent:H4}
The divisor covariance introduced in the companion paper has Fourier
symbol of order \(N\log N\) at quadratic scale, whereas the original von
Mangoldt jump has order \(\sqrt N(\log N)^2\); no extraction identity
connects them.  Under the literal Gamma--pole--threshold interpretation, a
fixed compact even core orthogonal to both \(q_L^+\) and the pole has
value \(-\kappa_N+O(1)=-4\lambda+o(\lambda)\), contradicting the proposed
\(-O(d_8)\) floor.
\end{entry}

\emph{Consequence.} We record this entry explicitly because the lemma in
question was ours.  Its viable successor is the complete coupled
prime--Gamma--pole complement inequality, i.e. the surplus.

\begin{entry}[\CTMD{} Qualitative positivity improvement is insufficient]
\label{ent:H5}
There is an explicit two-dimensional positivity-improving family with a
fixed negative ground state, a strictly positive near-radical model, and
overlap tending to zero at exactly the residual rate.
\end{entry}

\emph{Consequence.} Endpoint normalization is angle-blind and qualitative
positivity improvement cannot supply the denominator in the
anti-orthogonality quotient.  A viable positivity route must prove the
quantitative Harnack--tightness estimate for the actual ordinary-prime
operator.

\section{Three structural principles}
\label{sec:principles}

Every entry above is an instance of one of three principles.  We state
them because they are more useful than the individual entries: they
predict, without further computation, whether a proposed mechanism can
work.

\begin{proposition}[P1. Radical saturation]
The equality set of (2.7) contains the infinite-dimensional family
\(\{K^{(2j)}/K\}_{j\ge1}\), whose first member is real analytic and
unbounded.  Therefore no decomposition of
\(\EK-\frac12\Var_{\mu_K}\) into nonnegative summands may contain a
summand dominating a local edge square.
\end{proposition}

P1 governs C1, C2, C3, C4, C5, D1 and D3.  \emph{Diagnostic:} if a
proposed argument produces a nonnegative leftover supported on
displacements, it is already refuted.

\begin{proposition}[P2. Measure capacity]
The source carries positive singular mass on the lines
\(t\pm s=\log p^k\), while the polar pair law is absolutely continuous;
and the Gamma density alone lacks polar capacity on remote
fixed-displacement rectangles.  Therefore no endpoint-preserving positive
coupling, conditional expectation, martingale routing, or weight-preserving
charge from the source to the polar gradient exists.
\end{proposition}

P2 governs B1, B2, D1 and, through Jensen equality, the endpoint clause of
D2.  \emph{Diagnostic:} if a proposed argument routes source energy to
target energy through a positive kernel that respects endpoints, it is
already refuted.

\begin{proposition}[P3. Counterfactual invariance]
Under the existence of an off-line orbit, all of \(\Lambda(p^k)>0\),
\(K>0\), \(\mathcal J_{\log p^k}\ge0\) and \((hq)*K=0\) persist, while the
strict surplus fails; and the surplus itself equals the completed form at
the stationary residual.  Therefore no implication insensitive to the
location of the mean-periodic frequencies can prove it.
\end{proposition}

P3 governs A1--A3, E1--E7, F1--F6, G1--G5, H1--H5, and ultimately every
entry.  \emph{Diagnostic:} if a proposed argument would remain valid
verbatim in a system with an off-line zero, it is already refuted.  This
is the cheapest available audit and should be applied first.

\section{The admissible corridor}
\label{sec:corridor}

Negative results are only useful if they leave something.  We state what
the three principles jointly permit.

\begin{enumerate}
\item The decomposition must be \emph{signed}, with genuine cancellation
 between distinct prime channels (P1).  Individual squares are forbidden;
 signed differences of squares are not.
\item The transfer must be \emph{nonlocal} in the displacement variable,
 and must not preserve endpoints (P2).  It may be row dependent.
\item It must be applied \emph{after} the complete mean-periodic
 anti-short, and must use the actual prime-power locations \emph{before}
 any norm is taken (P3).  A mechanism that takes norms atom by atom has
 already discarded the placement information which is the only thing that
 distinguishes \(\zeta\) from the countermodels of D2, D3, E2, E3 and E4.
\end{enumerate}

Concretely, the surviving target admits two equivalent exact forms.  The
first is the displacement form
\[
 \boxed{
 \mathfrak P_{\rm PNT}(q)+\mathfrak b_{w_\Gamma}(q)
 +2\int_\R(K*K)K|q|^2\dd x\ \ge\ 0 ,}
\tag{13.1}
\]
with the sufficient uniform-margin version (3.13).  The second is the
relative operator form
\[
 \boxed{\mathcal C^\sharp\mathcal C\ \preceq\ \widetilde A ,}
\tag{13.2}
\]
equivalently \(\|H_C\|\le1\) for the minimal-norm compact factor of
Entry~G2.  By Theorem~\ref{thm:absorption} the Hermitian connection power
in (13.1) has already been absorbed with margin \(499/2000\); what remains
is exactly the common outgoing and incoming ordinary-prime power, which by
Entry~C5 may not be split, and by Entry~E7 may not be estimated on a
relative interval.

\begin{remark}[Honest accounting]
By Theorem~\ref{thm:antisurplus} both (13.1) and (13.2) are equivalent to
the Riemann Hypothesis, not weaker than it.  The corridor described here
is therefore a description of where a proof must live, not evidence that
one exists.  We state this plainly because the recurrence of the same
threshold constant across a dozen independent reformulations
(Entries G1, G2) is easily mistaken for convergence toward a proof, when
it is in fact a rigidity phenomenon: the constant is invariant precisely
because the problem has not moved.
\end{remark}

\section{Summary table}

\begin{center}
\small
\begin{longtable}{>{\raggedright\arraybackslash}p{0.055\textwidth}
                  >{\raggedright\arraybackslash}p{0.40\textwidth}
                  >{\raggedright\arraybackslash}p{0.115\textwidth}
                  >{\raggedright\arraybackslash}p{0.30\textwidth}}
\toprule
Entry & Mechanism excluded & Type & Principle / witness\\
\midrule
\endhead
A1 & Perron--Frobenius, Harnack, pointwise cone & \NOGO & Beurling--Deny fails; \(e^{d/2}\) vs \(e^{-d/2}\)\\
A2 & Variation-diminishing ordering & \NOGO & Gamma density log-convex, not \(\mathrm{TP}_2\)\\
A3 & Stochastic monotonicity, invariant cone & \NOGO & \(p=2\) tail-rate jump at \(t_0=\frac14\log2\)\\
B1 & Conditional expectation, first chaos & \NOGO & P2 capacity\\
B2 & Positive martingale routing & \NOGO & P2 capacity; two-state fibre \(2<S<3\)\\
B3 & Passing to the radical complement & \NOGO & Shorting is an equivalence\\
C1 & Canonical paths, Efron--Stein, Cheeger, local SOS & \NOGO & P1; \(K''/K\asymp e^{4x}\)\\
C2 & Local Hodge/Korn completion & \NOGO & P1; incidence cycles\\
C3 & Boundary flux at infinity & \NOGO & Conductance \(O(e^{-ce^R})\)\\
C4 & Gaussian/Volterra primitive squares & \NOGO & \(\eta<0\) on \((4/3,2)\)\\
C5 & Centered-path, local Stieltjes & \NOGO & \(w_*<0\) on \([\log\frac98,\log2)\); \(N^2\) loss\\
D1 & Cauchy--Binet charging & \NOGO & Circular; P2 for universal charges\\
D2 & Effective resistance, matrix-tree & \CTMD & Four-state rational model\\
D3 & Source algebra, inertia, summability & \CTMD & Literal-theta saturation model\\
D4 & Observation-only Gram ratio & \SCAL & Misses deficit by \(10^5\) at \(M=12\)\\
E1 & Finite-head Cholesky/Pick telescope & \NOGO & Head negative on every radical direction\\
E2 & Discrete levels + full-rank sensors & \CTMD & Three-state model, eigenvalue \(1/4\)\\
E3 & Tail energy exceeding deficit & \CTMD & Two-mode adaptation-loss model\\
E4 & Stieltjes positivity + summability & \CTMD & Exact equality saturation\\
E5 & Arbitrary-vector midpoint frame floor & \NOGO & Response--inertia bound\\
E6 & Local Christoffel observability & \NOGO & Deficit \(=\) full omitted energy\\
E7 & Relative-interval prime sampling & \SCAL & Loses \(e^{-2\pi\delta X}\)\\
F1 & Birman--Schwinger localization & \NOGO & Term-for-term equivalent\\
F2 & Mourre, dilation commutator & \NOGO & Odd/unbounded; vanishes on bound rows\\
F3 & Heat moments, operator-monotone maps & \NOGO & Equivalent to the surplus\\
F4 & Positive source atoms, heat moments & \NOGO & Compatible with \(\alpha<1/2\)\\
F5 & Finite-dimensional reduction as \(\delta\downarrow0\) & \SCAL & \(\Omega\gtrsim e^{C_R/\delta}\)\\
F6 & Central tightness from mean periodicity & \NOGO & \(T\log T\) vs \(O(T)\) zeros\\
G1 & Changing the signed kernel & \NOGO & Universal threshold \((2\inf\sigma A)^{-1/2}\)\\
G2 & Changing the realization & \NOGO & Minimal-norm compact factor\\
G3 & Multiplying passive local cells & \NOGO & \(K_L\) periodic falsifier\\
G4 & Nonlinear amplification & \NOGO & Leaves divisor; \(\zeta(s)/\zeta(s+\eps)\)\\
G5 & Bar homotopy, M\"obius inversion & \NOGO & \(\eps^{-m}\) gain, nonzero \(H_0\)\\
H1 & Fixed derivative models & \NOGO & Common factor \(\XiR\)\\
H2 & Many-vector co-Poisson frames & \NOGO & Frame identity; \(O(\lambda^2)\) vs \(O(Le^{C\lambda})\)\\
H3 & Unquantified cyclicity & \NOGO & Cost \(\gg\lambda^{2(1-\beta)}\)\\
H4 & Divisor-covariance extraction lemma & \NOGO & False as stated; \(-4\lambda+o(\lambda)\)\\
H5 & Qualitative positivity improvement & \CTMD & Two-dimensional family\\
\bottomrule
\end{longtable}
\end{center}

\section{Conclusion}

The companion paper reduces the Riemann Hypothesis, along the
Connes--CCM--Suzuki route, to the strict physical surplus.  This paper
establishes three things about that surplus.

First, it is not a lemma.  By Theorem~\ref{thm:antisurplus} it equals the
completed Weil form evaluated at the stationary residual, so a proof of it
is a proof of the Riemann Hypothesis and nothing less.  The practical form
of this observation is Remark~\ref{rem:notlemma}: audit any candidate
argument at the first line where (4.1) receives a positive sign.

Second, the space of mechanisms available to prove it is much smaller than
it appears.  Thirty-one entries in eight families are excluded, and the
exclusions are not independent: they follow from three principles ---
radical saturation, measure capacity, and counterfactual invariance ---
each of which can be checked against a proposed argument in a few lines.

Third, the missing constant is rigid.  By Entries G1 and G2, every exact
signed kernel and every realization of the same physical bank yields the
identical contraction threshold.  The repeated appearance of that constant
across independent reformulations is therefore not evidence of progress
toward it; it is a theorem that the problem is invariant under
reformulation.

What survives is narrow and precisely stated: a signed, globally nonlocal,
non-endpoint-preserving contraction applied after the complete
mean-periodic anti-short, using the actual prime-power placements before
any norm is taken, in either of the equivalent forms (13.1) or (13.2).
The unconditional gain of Theorem~\ref{thm:absorption} removes the
Hermitian connection power from that obligation with the explicit margin
\(499/2000\), leaving the common outgoing and incoming ordinary-prime
power as the sole remaining supply.

We close with the observation that motivated this paper.  A programme of
this kind produces far more refutations than theorems, and the refutations
are the part whose value does not depend on the conjecture ever being
settled.  We would have saved considerable time had a catalogue of this
kind existed when we began.

\begin{thebibliography}{99}

\bibitem{Weil1952}
A.~Weil,
\emph{Sur les ``formules explicites'' de la theorie des nombres premiers},
Comm. Sem. Math. Univ. Lund (1952), 252--265.

\bibitem{Connes1999}
A.~Connes,
\emph{Trace formula in noncommutative geometry and the zeros of the
Riemann zeta function},
Selecta Math. (N.S.) \textbf{5} (1999), 29--106.

\bibitem{Connes2026}
A.~Connes,
\emph{The Riemann Hypothesis: Past, Present and a Letter Through Time},
arXiv:2602.04022 (2026).

\bibitem{CCM2025}
A.~Connes, C.~Consani, and H.~Moscovici,
\emph{Zeta spectral triples},
arXiv:2511.22755 (2025).

\bibitem{CCJacobian}
A.~Connes and C.~Consani,
\emph{On the Jacobian of \(\overline{\operatorname{Spec}\mathbb Z}\)},
arXiv:2602.15941 (2026).

\bibitem{CCAbsolute}
A.~Connes and C.~Consani,
\emph{On the Absolute Geometry of \(\operatorname{Spec}\mathbb Z\)},
arXiv:2606.06604 (2026).

\bibitem{Suzuki2025}
M.~Suzuki,
\emph{On the Hilbert space derived from the Weil distribution},
Canadian Journal of Mathematics, FirstView (2025).

\bibitem{SuzukiScrew}
M.~Suzuki,
\emph{Weil's quadratic form via the screw function},
arXiv:2606.09096 (2026).

\bibitem{BeurlingDeny}
A.~Beurling and J.~Deny,
\emph{Dirichlet spaces},
Proc. Nat. Acad. Sci. U.S.A. \textbf{45} (1959), 208--215.

\bibitem{FukushimaOshimaTakeda}
M.~Fukushima, Y.~Oshima, and M.~Takeda,
\emph{Dirichlet Forms and Symmetric Markov Processes},
2nd ed., de Gruyter, 2011.

\bibitem{BirmanSolomyak}
M.~Sh. Birman and M.~Z. Solomyak,
\emph{Spectral Theory of Self-Adjoint Operators in Hilbert Space},
Reidel, 1987.

\bibitem{Douglas}
R.~G. Douglas,
\emph{On majorization, factorization, and range inclusion of operators on
Hilbert space},
Proc. Amer. Math. Soc. \textbf{17} (1966), 413--415.

\bibitem{Haynsworth}
E.~V. Haynsworth,
\emph{Determination of the inertia of a partitioned Hermitian matrix},
Linear Algebra Appl. \textbf{1} (1968), 73--81.

\bibitem{Karlin}
S.~Karlin,
\emph{Total Positivity, Vol. I},
Stanford University Press, 1968.

\bibitem{TrejoPizzo40}
D.~A. Trejo Pizzo,
\emph{Branch Selection in the Connes--CCM Route to RH: Weil Positivity,
Prolate Zeta-Cycles, and Cofinal Prime--Gamma Inertia Matching},
companion paper (2026).

\end{thebibliography}

\end{document}
